\documentclass[12pt]{article}
\usepackage[utf8]{inputenc}
\usepackage[T1]{fontenc}
\usepackage{hyperref}
\title{Setaman Language Resources for K-12 Curriculum Development}
\author{Compiled by Claude Sonnet 5.5 (claude-sonnet-5-5) for Andr\'as Kornai}
\date{}
\begin{document}
\maketitle

\section{Introduction}
Setaman (ISO 639-3 \texttt{stm}) is a Trans-New Guinea (Ok branch) language of the Western Province of Papua New Guinea. It is spoken around the head of the Murray River in villages such as Seltamin, Ngisokabip and Baktamin. About 500 speakers are reported. Some sources treat it as a dialect of Faiwol. Very little published material exists.

\section{Primary Educational Resources}
No primers or textbooks were found. A sociolinguistic survey reports that there are no schools in the language area. Children who go to school leave the area and are taught in Faiwol, Tok Pisin or English.
\begin{itemize}
\item A Sociolinguistic Survey of Setaman [stm]: \url{https://www.endangeredlanguages.com/node/97632}
\end{itemize}

\section{Linguistic and Reference Materials}
\begin{itemize}
\item Wikipedia, ``Faiwol language'' (Setaman is listed as a dialect or related variety): \url{https://en.wikipedia.org/wiki/Faiwol_language}
\item Numeral systems database entry for Setaman (body-part counting system documented by Barth, 1975): \url{https://lingweb.eva.mpg.de/channumerals/Setaman.htm}
\end{itemize}

\section{Religious and Community Resources}
\begin{itemize}
\item Joshua Project profile. It notes that no Scripture exists in Setaman and that the area has no health facilities or schools: \url{https://joshuaproject.net/people_groups/14875}
\end{itemize}

\section{Digital and Audio Resources}
No audio, video or software was found.

\section{Curriculum Development Context}
Setaman is spoken as a first language by the whole community and speakers view it positively. Code-switching is reported to be rare. There is no local school and no written materials, so any curriculum would start from scratch. Faiwol, Tok Pisin and English are the languages of schooling for children outside the area. Little else exists.

\end{document}